%
% mail-refile-confirm-early.tex  (FROZEN NEGATIVE CONTROL: do not edit)
%
% Confirm has no first w? in shown guard: the client confirms before it
% prints.  Expected: MailInv violated.
%
% Z Specification for a partial read on a consuming relay: the refile of the
% remainder, and the confirmation of the batch (DES-095, DES-099).
%
% Companion to session-model.tex and nats-relay.tex.  session-model.tex keeps
% the unread mail of a session as a count ($durable$, $shown$); this model
% needs the messages themselves, because the properties are about which
% message is where.  It is kept apart because those two specifications
% already have large state spaces.
%
% Type-check: fuzz -t mail-refile.tex
% Model-check: see the end of this file.
%

\documentclass[a4paper,10pt,fleqn]{article}
\usepackage[margin=1in]{geometry}
\usepackage{fuzz}
\usepackage[colorlinks=true,linkcolor=blue,citecolor=blue,urlcolor=blue]{hyperref}

\begin{document}

\title{Biff Mail Refile and Confirmation: A Z Specification}
\author{Punt Labs}
\date{November 2026}
\maketitle

%%----------------------------------------------------------------------
\section{Introduction}
%%----------------------------------------------------------------------

The NATS relay \emph{consumes} mail when it is fetched: a fetch acknowledges
the messages it returns, and the work queue deletes what is acknowledged.
After a fetch the mail exists only in the memory of whoever fetched it.

The daemon's \texttt{read} fetches a bounded batch, fits as much of it into
the reply as the read budget allows, and is left with a remainder.  The
remainder used to stay in the daemon's outbox, out of reach of everything else
that drains a session: ending the session, or a read from another process,
found the relay empty.  A test on a real relay found 0 of 26 messages
reachable after a partial read.  The remainder is now filed back into the
relay under an id derived from its own (\texttt{Message.refiled}), so that
filing it twice cannot make two copies.

The reply carries the batch.  The batch stays in the outbox until the client
confirms it, and the client confirms only after it has printed it
(\texttt{request\_held}, \texttt{HeldReply.confirm}).  A confirmation sent
before the print let a client that died in between mark as read a message it
had never shown.

This specification states what is claimed and what is not.  The claims hold
when nothing fails; each failure that voids one is an explicit operation.

\begin{enumerate}
  \item \textbf{No remainder is lost.} After any partial read, every message
    the batch does not carry is in the relay or, if the relay refused it,
    in the outbox of a session that has not ended.
  \item \textbf{Nothing is duplicated.} A message is held by at most one of
    the relay and the outbox, and the relay holds at most one copy of it.
  \item \textbf{\texttt{more} counts each remaining message once.}
  \item \textbf{A message is marked read only after it was shown.}
  \item \textbf{Nothing vanishes unaccounted.} Every issued message is shown,
    drained, held, or recorded as lost by a named event.
\end{enumerate}

%%----------------------------------------------------------------------
\section{Basic Types}
%%----------------------------------------------------------------------

\begin{zed}
  ZBOOL ::= ztrue | zfalse
\end{zed}

\begin{zed}
  ReadPhase ::= rdIdle | rdFetched
\end{zed}

\begin{axdef}
  fetchMax : \nat \\
  batchMax : \nat
  \where
  fetchMax = 2 \\
  batchMax = 1
\end{axdef}

$fetchMax$ is the bound of one consuming fetch; $batchMax$ is the number of
messages the read budget fits into one reply.  Both are small so that
ProB can enumerate them, and $batchMax < fetchMax$ is what makes a read
partial.

\begin{zed}
  MAIL == 1 \upto 2 \\
  GEN == 0 \upto 1 \\
  WIRE == MAIL \cross GEN
\end{zed}

$MAIL$ is a message as the sender knows it.  $WIRE$ is a message as the relay
knows it: the message and the generation of its id.  A refile gives the
message the next id in its chain; the id is \texttt{uuid5} of the previous
one, so the chain never meets itself.  The model counts the generations
modulo two, which is sound because a message has one live copy (property 2),
and the one pair that would collide is exactly the case the derivation is
there for: a copy filed twice from the same predecessor.

\begin{axdef}
  refile : WIRE \fun WIRE
  \where
  refile = \{ m : MAIL; g : GEN @ (m, g) \mapsto (m, (g + 1) \mod 2) \}
\end{axdef}

\begin{axdef}
  mailOf : \power WIRE \fun \power MAIL
  \where
  \forall s : \power WIRE @ mailOf~s = \{ w : s @ first~w \}
\end{axdef}

%%----------------------------------------------------------------------
\section{State}
%%----------------------------------------------------------------------

\begin{schema}{Store}
  relay : \power WIRE \\
  held : \power WIRE \\
  sent : \power WIRE \\
  phase : ReadPhase
\end{schema}

$relay$ is the session's inbox in the relay, shared with every process that
drains it.  $held$ is the daemon's outbox: process memory, lost with the
process.  $sent$ is the batch of the reply the client has not yet confirmed;
a message in $sent$ is held.  $phase$ is $rdFetched$ between the fetch of a
read and its reply.

\begin{schema}{Client}
  shown : \power MAIL \\
  done : \power MAIL \\
  dupShown : ZBOOL
\end{schema}

$shown$ is what the client has printed, $done$ what has been marked read, and
$dupShown$ records that a reply offered the client a message it had already
printed.

\begin{schema}{Ledger}
  issued : \power MAIL \\
  drained : \power MAIL \\
  lostBatch : \power MAIL \\
  lostRest : \power MAIL
\end{schema}

A ghost.  $issued$ is every message ever filed; $drained$ is what an ending
session moved into its owner's mailbox, where it is reachable.  $lostBatch$
and $lostRest$ are what a failure took from the outbox without the client
having been shown it: the batch of an unconfirmed reply, and the rest.

\begin{schema}{Flags}
  faulted : ZBOOL \\
  moreOut : 0 \upto 2 \\
  moreTruth : 0 \upto 2
\end{schema}

$faulted$ is set by every operation that is a failure and not a protocol
step: a daemon crash, a relay that refuses a refile, a read cancelled
half-way, a session that ends during a read.  $moreOut$ is the \texttt{more}
the last reply carried, computed as the code computes it; $moreTruth$ is the
number of messages that reply left, computed from the ledger.

\begin{schema}{MailState}
  Store \\
  Client \\
  Ledger \\
  Flags
\end{schema}

\begin{schema}{MailInv}
  MailState
  \where
  sent \subseteq held \\
  done \subseteq shown \\
  shown \subseteq issued \\
  issued \subseteq shown \cup drained \cup mailOf~relay \cup mailOf~held \cup
    lostBatch \cup lostRest \\
  faulted = zfalse \implies lostRest = \emptyset \\
  faulted = zfalse \implies \\
  \t1 ( \forall m : MAIL @ \# \{ w : relay \cup held | first~w = m \} \leq 1 ) \\
  faulted = zfalse \implies moreOut = moreTruth
\end{schema}

The five properties, in the order of the introduction (the first is the
fifth line, the last the second).

\begin{schema}{Invariant}
  MailInv
\end{schema}

\begin{schema}{Init}
  MailState'
  \where
  relay' = \emptyset \\
  held' = \emptyset \\
  sent' = \emptyset \\
  phase' = rdIdle \\
  shown' = \emptyset \\
  done' = \emptyset \\
  dupShown' = zfalse \\
  issued' = \emptyset \\
  drained' = \emptyset \\
  lostBatch' = \emptyset \\
  lostRest' = \emptyset \\
  faulted' = zfalse \\
  moreOut' = 0 \\
  moreTruth' = 0
\end{schema}

%%----------------------------------------------------------------------
\section{Operations}
%%----------------------------------------------------------------------

\begin{schema}{Arrive}
  \Delta Store \\
  \Xi Client \\
  \Delta Ledger \\
  \Xi Flags \\
  m? : MAIL
  \where
  m? \notin issued \\
  relay' = relay \cup \{ (m?, 0) \} \\
  held' = held \\
  sent' = sent \\
  phase' = phase \\
  issued' = issued \cup \{ m? \} \\
  drained' = drained \\
  lostBatch' = lostBatch \\
  lostRest' = lostRest
\end{schema}

A message is filed in the session's inbox.

\begin{schema}{ReadFetch}
  \Delta Store \\
  \Xi Client \\
  \Xi Ledger \\
  \Xi Flags \\
  fetched? : \power WIRE
  \where
  phase = rdIdle \\
  fetched? \subseteq relay \\
  \# fetched? \leq fetchMax \\
  \# relay \leq fetchMax \implies fetched? = relay \\
  relay' = relay \setminus fetched? \\
  held' = held \cup fetched? \\
  sent' = sent \\
  phase' = rdFetched
\end{schema}

The consuming fetch: what it returns leaves the relay and is held.  A
bounded fetch takes the whole inbox when the inbox fits.

\begin{schema}{ReadReply}
  \Delta Store \\
  \Delta Client \\
  \Xi Ledger \\
  \Delta Flags \\
  batch? : \power WIRE
  \where
  phase = rdFetched \\
  batch? \subseteq held \\
  shown' = shown \\
  done' = done \\
  dupShown' = ( \IF mailOf~batch? \cap shown \neq \emptyset \THEN ztrue \ELSE dupShown ) \\
  \# batch? \leq batchMax \\
  \# held \leq batchMax \implies batch? = held \\
  \# held > batchMax \implies \# batch? = batchMax \\
  relay' = relay \cup refile \limg held \setminus batch? \rimg \\
  held' = batch? \\
  sent' = batch? \\
  phase' = rdIdle \\
  faulted' = faulted \\
  moreOut' = \# relay' \\
  moreTruth' = \# ( mailOf~( relay' \cup ( held' \setminus batch? ) ) \setminus
    mailOf~batch? )
\end{schema}

The reply and the refile.  The batch is the part of the outbox the budget
fits.  The remainder is filed back, each message under its derived id, and
released from the outbox; the batch stays, unconfirmed.  The reply says
$moreOut' = \# relay'$: \texttt{left = len(pending) - len(batch) - filed +
unfetched}, where $filed$ is the remainder and $unfetched$ the relay's own
count after the refile, so the remainder is subtracted from the one count and
added by the other.  $moreTruth'$ is the number of \emph{distinct messages} the
reply leaves: those held anywhere for a later read, less the ones this
reply carries.  The two agree when no message has two copies; that is the
sense in which \texttt{more} counts each remaining message once.

\begin{schema}{ReadReplyRefused}
  \Delta Store \\
  \Delta Client \\
  \Xi Ledger \\
  \Delta Flags \\
  batch? : \power WIRE
  \where
  phase = rdFetched \\
  batch? \subseteq held \\
  shown' = shown \\
  done' = done \\
  dupShown' = ( \IF mailOf~batch? \cap shown \neq \emptyset \THEN ztrue \ELSE dupShown ) \\
  \# batch? \leq batchMax \\
  \# held \leq batchMax \implies batch? = held \\
  \# held > batchMax \implies \# batch? = batchMax \\
  relay' = relay \\
  held' = held \\
  sent' = batch? \\
  phase' = rdIdle \\
  faulted' = ztrue \\
  moreOut' = moreOut \\
  moreTruth' = moreTruth
\end{schema}

The relay refused the refile (\texttt{RELAY\_ERRORS}, logged).  The remainder
stays in the outbox and is offered again by the next read.

\begin{schema}{ReadCancelled}
  \Delta Store \\
  \Xi Client \\
  \Xi Ledger \\
  \Delta Flags \\
  rest? : \power WIRE
  \where
  phase = rdFetched \\
  rest? \subseteq held \\
  rest? \neq \emptyset \\
  relay' = relay \cup refile \limg rest? \rimg \\
  held' = held \\
  sent' = sent \\
  phase' = rdIdle \\
  faulted' = ztrue \\
  moreOut' = moreOut \\
  moreTruth' = moreTruth
\end{schema}

The read is cancelled after some of its refiles landed and before the
outbox released them.  The connection's closing cancels its running requests
(\texttt{Connection.\_serve}, \texttt{task.cancel()}), and the release is one
call after the whole loop, so this is reachable.  Each such message is in the
relay and in the outbox.

\begin{schema}{ClientShow}
  \Xi Store \\
  \Delta Client \\
  \Xi Ledger \\
  \Xi Flags \\
  w? : WIRE
  \where
  w? \in sent \\
  shown' = shown \cup \{ first~w? \} \\
  done' = done \\
  dupShown' = dupShown
\end{schema}

The client prints a message of the reply it holds.

\begin{schema}{Confirm}
  \Delta Store \\
  \Delta Client \\
  \Xi Ledger \\
  \Xi Flags \\
  w? : WIRE
  \where
  w? \in sent \\
  held' = held \setminus \{ w? \} \\
  sent' = sent \setminus \{ w? \} \\
  relay' = relay \\
  phase' = phase \\
  shown' = shown \\
  done' = done \cup \{ first~w? \} \\
  dupShown' = dupShown
\end{schema}

The client confirms a message, and the daemon marks it read and forgets it.
The guard $first~w? \in shown$ is the ordering of \texttt{Answer.deliver}:
print, flush, then confirm.

\begin{schema}{SessionEnd}
  \Delta Store \\
  \Xi Client \\
  \Delta Ledger \\
  \Xi Flags
  \where
  phase = rdIdle \\
  relay' = \emptyset \\
  held' = \emptyset \\
  sent' = \emptyset \\
  phase' = rdIdle \\
  issued' = issued \\
  drained' = drained \cup mailOf~relay \\
  lostBatch' = lostBatch \cup ( mailOf~sent \setminus shown ) \\
  lostRest' = lostRest \cup ( mailOf~( held \setminus sent ) \setminus shown )
\end{schema}

The session ends.  The drain moves what the relay holds into the owner's
mailbox, and the session's record, with its outbox, is discarded.

\begin{schema}{SessionEndDuringRead}
  \Delta Store \\
  \Xi Client \\
  \Delta Ledger \\
  \Delta Flags
  \where
  phase = rdFetched \\
  relay' = \emptyset \\
  held' = \emptyset \\
  sent' = \emptyset \\
  phase' = rdIdle \\
  issued' = issued \\
  drained' = drained \cup mailOf~relay \\
  lostBatch' = lostBatch \cup ( mailOf~sent \setminus shown ) \\
  lostRest' = lostRest \cup ( mailOf~( held \setminus sent ) \setminus shown ) \\
  faulted' = ztrue \\
  moreOut' = moreOut \\
  moreTruth' = moreTruth
\end{schema}

The session ends between a read's fetch and its reply.

\begin{schema}{DaemonCrash}
  \Delta Store \\
  \Xi Client \\
  \Delta Ledger \\
  \Delta Flags
  \where
  relay' = relay \\
  held' = \emptyset \\
  sent' = \emptyset \\
  phase' = rdIdle \\
  issued' = issued \\
  drained' = drained \\
  lostBatch' = lostBatch \cup ( mailOf~sent \setminus shown ) \\
  lostRest' = lostRest \cup ( mailOf~( held \setminus sent ) \setminus shown ) \\
  faulted' = ztrue \\
  moreOut' = moreOut \\
  moreTruth' = moreTruth
\end{schema}

The daemon dies.  The relay keeps what it holds; the outbox is gone.

\begin{schema}{Quiesce}
  \Xi MailState
\end{schema}

The stutter that stands for ``nothing more happens'': the bound on
$MAIL$ ends the history of arrivals, and without $Quiesce$ ProB would
report the end of the bounded history as a deadlock.

%%----------------------------------------------------------------------
\section{Model Checking}
%%----------------------------------------------------------------------

$MailInv$ is the state invariant; probcli reports it for the shipped
operations in the transcript at the end of this file.

Three things the specification does \emph{not} promise, each reachable and
each stated as a reachability formula so that it is not an omission:

\begin{enumerate}
  \item \textbf{A duplicate delivery is possible.}  $EF\;(dupShown = ztrue)$
    holds: a client that prints, dies before it confirms, and reads again is
    shown the message again, because the batch is still in the outbox and the
    next read offers it.  Delivery is at least once.  The same read also
    refiles a message that an earlier, unconfirmed reply already showed.
  \item \textbf{A batch is lost if the daemon dies or the session ends
    before the client has been shown it.}  $EF\;(lostBatch \neq \emptyset)$
    holds, along the trace $Arrive \to ReadFetch \to ReadReply \to
    SessionEnd$ with no $ClientShow$: the batch of an unconfirmed reply is
    held only in the outbox, and neither the drain nor a restart reads the
    outbox.  A client that dies between the reply and its print keeps the
    message unread \emph{only while the session and the daemon live}.
  \item \textbf{A fault voids the counts.}  A cancelled read leaves a copy in
    the relay and in the outbox, so the relay's next reply counts it twice.
\end{enumerate}

Two frozen variants are the negative controls.
\texttt{mail-refile-drop.tex} has a $ReadReply$ that files nothing back
(the code before the refile): $MailInv$ is violated, with the trace $Arrive
\to Arrive \to ReadFetch \to ReadReply \to SessionEnd$, the remainder lost
with no fault.  \texttt{mail-refile-confirm-early.tex} has a $Confirm$ with no
$first~w? \in shown$ guard (the confirmation sent before the print):
$MailInv$ is violated by a message marked read and never shown.

%%----------------------------------------------------------------------
%% Reproducing the checks.  Run from the repo root; one probcli at a time.
%%
%%   PROBCLI="${PROBCLI:-$HOME/Applications/ProB/probcli}"
%%   P="-p DEFAULT_SETSIZE 1 -p MAXINT 6 -p MAX_INITIALISATIONS 20 \
%%      -p MAX_OPERATIONS 500 -p TIME_OUT 600000"
%%   "$PROBCLI" docs/specs/mail-refile.tex $P -model_check
%%   "$PROBCLI" docs/specs/mail-refile.tex $P -model_check -ctlformula "EF {dupShown=ztrue}"
%%   "$PROBCLI" docs/specs/mail-refile.tex $P -model_check -ctlformula "EF {lostBatch/={}}"
%%   "$PROBCLI" docs/specs/mail-refile.tex $P -model_check -ctlformula "EF {moreOut>0}"
%%
%% Negative controls, each of which must report an invariant violation:
%%
%%   "$PROBCLI" docs/specs/mail-refile-drop.tex $P -model_check
%%   "$PROBCLI" docs/specs/mail-refile-confirm-early.tex $P -model_check
%%----------------------------------------------------------------------

\end{document}
